\subsection{Interpretation of the Final Factorial Experiment}

The most important finding of the final COVA-3D experiment is negative but
informative. None of the six sparse-supervision configurations achieved a
non-zero lesion recall at the prospectively specified low-false-positive
operating points. This result was consistent across three independently trained
seeds and remained unchanged when the allowed false-positive burden was
increased from 0.5 to 4 false positives per patient.

The failure of the primary endpoint should not be interpreted as an absence of
all lesion-related information in the model outputs. The secondary analysis
showed substantially higher any-overlap lesion recall than geometrically stricter
IoU-based lesion recall. This indicates that predicted regions frequently reached
true lesions but did not form sufficiently coherent lesion instances. The very
large number of connected prediction components and high false-positive volume
provide a complementary explanation for the failure of the FROC operating
points: lesion-contact information was present, but it was accompanied by severe
topological fragmentation and background activation.

The primary factorial hypothesis was therefore not supported. Under the tested
baseline architecture, training schedule, fixed sparse annotation budgets, and
frozen low-FP evaluation protocol, changing lesion-instance coverage or
within-lesion scribble geometry was insufficient to produce reliable
low-false-positive lesion recovery. In particular, the two prospectively defined
coverage-by-geometry interactions were exactly zero under the primary endpoint.

The descriptive secondary measures nevertheless suggest that annotation geometry
can alter the character of the prediction error. Coherent supervision produced
comparatively high lesion-contact recall, while fragmented and dispersed
conditions showed different overlap, component-count, and false-positive
profiles. Because these analyses were not accompanied by prospectively defined
confirmatory secondary tests, they should be interpreted as mechanistic
observations rather than evidence of superiority between factorial cells.

These findings identify a practical limitation of the present sparse-supervision
baseline. Future method development should focus on controlling component
proliferation, suppressing false-positive regions, and improving lesion-instance
coherence rather than merely reallocating the same sparse annotation budget.
Such follow-up development should be treated as a new prospective experiment and
must not alter or replace the frozen results reported here.
